\section{Methodology}
\label{sec:method}

\subsection{State Representation}
\label{sec:state}

\textbf{Position.}
The world-frame position $\vp_w(t) \in \mathbb{R}^3$ is a quintic
(order-6) uniform B-spline with knot spacing $\Delta t_p = 5$\,ms
(slow racing; per-regime grids in \autoref{tab:config}).
Velocity $\dot{\vp}_w$ and acceleration $\ddot{\vp}_w$ are
obtained as analytical derivatives of the spline.

\textbf{Orientation.}
We use the \emph{cumulative} SO(3) B-spline
formulation~\cite{sommer2020efficient}:
\begin{equation}
  \vR(t) = \vR_{\mathrm{base}}[k{-}3]
           \prod_{j=k-3}^{k} \Exp\!\bigl(\tilde{B}_j(t)\,\vOm_j\bigr),
  \label{eq:ori_spline}
\end{equation}
where $k$ is the active knot index for time $t$,
$\vOm_j \in \so$ are incremental rotation control points
(standard $\Exp$/$\Log$ map convention),
$\tilde{B}_j(t)$ are cumulative cubic basis functions, and
$\vR_{\mathrm{base}}$ is the left anchor rotation.  Knot spacing
is $\Delta t_o = 8$\,ms (slow racing; see \autoref{tab:config}).
Body-frame angular velocity $\vom_b(t)$
is obtained as a closed-form derivative of~\eqref{eq:ori_spline}
directly from \texttt{evaluate\_lie()} without numerical
differentiation.

\textbf{Biases and extrinsics.}
A constant 6-DoF \ac{imu} bias $[\vb_a, \vb_g]$ is estimated jointly.
The extrinsic pitch $\delta_\text{pitch}$ (1 DoF) is optimized
with a soft prior; roll and yaw are fixed (unobservable from
Doppler).  It is freed in the batch solve but locked at the measured
27.5\degree\ for the sliding window (\autoref{sec:backflips}).

\subsection{Measurement Models}
\label{sec:models}

\textbf{Radar Doppler with lever arm.}
The TI IWR6843 reports positive Doppler for a receding target.
With antenna offset $\mathbf{t}_{bs}$ from the body \ac{com}, the
predicted Doppler for a point at sensor-frame bearing $\vu_s$ is:
\begin{equation}
  v_D = -\vu_b^{\top}\!\left(
    \vR_{bw}\,\vv_w(t) + \vom_b(t) \times \mathbf{t}_{bs}
  \right),
  \label{eq:doppler}
\end{equation}
where $\vu_b = \vR_{bs}\,\vu_s$, $\vR_{bs}$ is the rotation from
radar sensor to body frame (extrinsic $[\text{roll}{=}180\degree,\,
\text{pitch}{=}27.5\degree,\,\text{yaw}{=}0\degree]$; measured,
\autoref{sec:results}), and
$\vR_{bw} = \vR_{wb}^{\top}$.  The leading minus makes the convention
self-consistent: a sensor translating toward a static point
($\vu_b^{\top}\vR_{bw}\vv_w > 0$) yields $v_D < 0$, an approaching rather
than receding target, as required.  The residual
$r_k = v_{D,\text{meas}} - v_{D,\text{pred}}$ uses Huber
loss with $\delta = 1.0$\,m/s, robust as defense-in-depth to any
residual Doppler noise on the surviving inlier returns; the
\SI{0.049}{\metre\per\second} Doppler resolution sets a lower bound but
does not drive the choice.

\textbf{\ac{ransac} ego-velocity prefilter (default front-end).}
A single-chip radar with only 2\,TX antennas has limited elevation
diversity, and a minority of returns carry an elevation-correlated
Doppler bias that a robust kernel only down-weights.  Following
\mbox{reve}~\cite{doer2020ekf} (Doer's radar ego-velocity estimator), we
therefore prepend a \mbox{reve}-style 3D
least-squares \ac{ransac} ego-velocity estimate to the per-frame Doppler set
(inlier threshold $0.15$\,m/s) and pass only the inlier returns to the
spline solver, hard-rejecting the elevation-biased minority before the
optimization.  The prefilter is seeded
(\texttt{numpy.default\_rng(0)}), so the full pipeline reproduces
bit-identically; frames with fewer than five returns bypass it
unchanged.  On those sparse frames, including most of the backflip
regime, the in-solver Huber loss is the sole radar outlier protection,
so the two stages are complementary, not redundant.  It runs once at
frame load, not per window, so it does not
enter the per-window timing budget (\autoref{sec:results}).
\autoref{sec:baselines} quantifies what it removes: on the public
ICINS flights it closes an order-of-magnitude vertical-drift gap, and on
our own flights it improves aggressive-racing live position by $20\%$ ($0.50\to0.40$\,m
live).  It is enabled by default for every result in this paper.

\textbf{Dynamics-adaptive measurement weighting.}
The central practical finding of this work is that a single
\emph{rate-adaptive trust allocation} across the three measurement
streams covers hover to \SI{10}{\radian\per\second} tumbling with one
configuration.  Each component was derived from a measured failure
diagnosis (\autoref{sec:backflips}), not tuned ad hoc:
\begin{itemize}
\item \emph{Gyroscope: trust more when rotating fast.}
  Per-sample weight
  $\lambda_g^{\mathrm{eff}} = \lambda_g\bigl(1 + (|\vz_g|/\omega_g)^p\bigr)$
  with $\lambda_g{=}4$, $\omega_g{=}4$\,rad/s, $p{=}4$, using the
  \emph{measured} rate (causal, no spline evaluation).  The quartic
  exponent keeps racing untouched ($\lambda^{\mathrm{eff}}\!\approx\!\lambda_g$
  below 3\,rad/s) and reaches $\lambda^{\mathrm{eff}}\!\approx\!160$--$500$
  during flips, where the gyroscope is the only clean sensor.  A
  quadratic ramp ($p{=}2$) stiffens the gyro already at 2--5\,rad/s
  (where radar is still informative) and degrades racing roll/yaw.
\item \emph{Radar: trust less when rotating fast.}
  A radar frame integrates over tens of milliseconds; at body rate
  $\omega$ the scene smears by $\omega\,T_{\mathrm{frame}}$
  ($\approx$17\degree\ at 10\,rad/s).  Per-frame noise inflation
  $\sigma_{\mathrm{eff}}^2 = \sigma_0^2\bigl(1+(|\vom|/\omega_r)^2\bigr)$,
  $\omega_r{=}4$\,rad/s, $|\vom|$ from the warm-start spline at
  problem-build time; continuous, no hard-gate tuning cliff.
\item \emph{Accelerometer: same form,} $\omega_a{=}8$\,rad/s:
  during flips the accelerometer carries thrust transients and
  vibration, its gravity information is nil mid-flip, and it otherwise
  distorts orientation through the $\vR$--$\ddot{\vp}$ coupling.
\item \emph{Per-point \ac{snr} weighting.}
  $w_i = \mathrm{clip}\bigl((I_i/\tilde{I})^{\alpha},\,0.25,\,4\bigr)$
  with frame-median intensity $\tilde I$ and $\alpha{=}1$
  (median-relative, so the global radar weight is unchanged).  High-\ac{snr}
  returns carry measurably better Doppler: on aggressive racing this
  alone improves live orientation from 3.31\degree\ to 2.88\degree\ at
  unchanged position.
\end{itemize}
The weighting trades are made explicit in
\autoref{sec:backflips}: on tumbling flight the radar/accelerometer
gates exchange velocity accuracy for orientation accuracy along a
measured Pareto curve; $(\omega_r,\omega_a){=}(4,8)$ is its knee.

\textbf{Flip-regime Doppler correction.}
Under sustained tumbling an additional attitude-correlated Doppler
systematic appears, which we absorb with a per-point term
$v_{\mathrm{corr}} = v_{\mathrm{meas}} - b\,u_{s,z}$ ($b{=}{-}1.5$\,m/s,
backflips only).  \autoref{sec:backflips} shows why $b$ must be
interpreted as a flight-regime error proxy rather than a physical
elevation bias of the 2-TX array.

\textbf{Accelerometer.}
\begin{equation}
  \mathbf{r}_a = \vz_a - \vR_{bw}\!\left(\ddot{\vp}_w - \vg_w\right) - \vb_a,
  \label{eq:accel}
\end{equation}
weighted $\lambda_a = 0.01$.

\textbf{Gyroscope.}
\begin{equation}
  \mathbf{r}_g = \vz_g - \vom_b(t) - \vb_g,
  \label{eq:gyro}
\end{equation}
weighted $\lambda_g = 4.0$ (L2 loss).

\textbf{Gravity direction.}
\begin{equation}
  \mathbf{r}_{\mathrm{tilt}} =
    \frac{\vz_a - \vb_a}{\|\vz_a - \vb_a\|} - \vR_{bw}^{\top}\hat{\vg},
  \label{eq:tilt}
\end{equation}
active when $\bigl|\,\|\vz_a{-}\vb_a\| - g\,\bigr| < 3\,\mathrm{m/s}^2$.
This provides roll/pitch anchoring during near-hover phases without
contaminating dynamic flight.

\textbf{Heading prior.}
Yaw is unobservable from Doppler alone (all yaw-equivalent
trajectories predict identical radial velocities under pure
rotation about gravity).  We supply a heading reference from a magnetometer (standard on
most platforms; here simulated by \ac{mocap} yaw).  The residual $r_\psi = \psi_\text{est} - \psi_\text{ref}$
is added at $\sim$100\,Hz, with $\lambda_\psi = 0.6$ by default and
raised to $10$ on the heading-critical bags (slow racing, backflips;
\autoref{tab:sw} footnote).
A weight sweep shows the split is a yaw-observability/position trade.
On slow racing, gentle dynamics leave yaw radar-unconstrained, so the
weak default ($\lambda_\psi{=}0.6$) lets yaw drift to $4.3\degree$
(orientation $2.0\to4.2\degree$) while a strong prior costs almost no
position.  On fast racing, the rich translational excitation already
holds yaw near $3.9\degree$ at $\lambda_\psi{=}0.6$, so raising it to
$10$ buys only $0.6\degree$ of yaw yet degrades settled position
($0.31\to0.43$\,m) through the orientation--velocity coupling; we
therefore keep it light.

\subsection{Regularization}
\label{sec:reg}

\textbf{Position: minimum snap.}
$E_{\text{snap}} = \lambda_s \int \|\vp^{(4)}(t)\|^2\,dt$
with $\lambda_s = 2{\times}10^{-5}$.

\textbf{Orientation: minimum angular acceleration.}
Rather than penalizing angular velocity increments (minimum $\omega$),
which opposes every banked turn and the entire backflip maneuver, we
penalize the \emph{second} finite difference on SO(3):
\begin{equation}
  \mathbf{r}_\alpha = \Log(\vR_{i-1}^{-1}\vR_i)
                    - \Log(\vR_i^{-1}\vR_{i+1}),
  \label{eq:min_alpha}
\end{equation}
which is zero for constant angular rate.
Swept $\lambda_\alpha \in \{0, 0.001, 0.01, 0.1, 1.0\}$; 0.1
was best across all bags (1.0 caused backflip divergence).

\subsection{Sensor-Only Initialization (P1--P3)}
\label{sec:init}

\textbf{P1: Orientation.}
A stationary segment before flight is detected automatically; the gyroscope
bias is the mean reading over this segment and the accelerometer bias a
gravity-aligned scalar correction, both \emph{sensor-only} and using no
external attitude reference.
Initial roll/pitch are then derived from the stationary accelerometer
(gravity direction).  Yaw is set to zero (gauge freedom; corrected
later by the heading prior).  The full orientation trajectory is
obtained by forward-integrating debiased gyroscope readings:
$\vR(t+\Delta t) = \vR(t)\,\Exp\!\bigl((\vz_g - \vb_g)\Delta t\bigr)$.

\textbf{P2: Position.}
Per radar frame, weighted least squares estimates the 3D ego-velocity
$\vv_{\text{wls}}$ from the \ac{ransac} inlier returns
(\autoref{sec:models}).  Doppler alias
unwrapping uses \ac{imu}-integrated world-frame velocity to select the
correct alias once at frame load (20.9\% of points aliased
in fast racing, 1.7\% in slow racing); the unwrapped Doppler is the
measurement consumed by \emph{both} this \ac{wls} estimate and the main
solver's radar factors (\autoref{eq:doppler}), with the solver's own
$v_{\max}$ check kept only as a residual safety net.
World-frame velocity is integrated with forward Euler to produce an
initial position trajectory.

\textbf{P3: Boundary priors.}
Position, velocity, orientation, and angular velocity boundary
priors are built from the sensor-only state at $t=0$.
$\lambda_{\text{bnd}} = 1000$ for position and velocity.

\subsection{Sliding Window with Schur Complement Marginalization}
\label{sec:sw}

\begin{figure}[t]
\centering
\scalebox{0.97}{%
\begin{tikzpicture}[
  font=\small,
  cp/.style={circle, draw, minimum size=6pt, inner sep=0pt},
  ]
  % --- shaded zone backgrounds (draw first, CPs on top) ---
  \fill[red!12]   (1.5,-0.22) rectangle (2.7, 0.22);
  \fill[green!12] (2.7,-0.22) rectangle (4.5, 0.22);

  % --- timeline ---
  \draw[->, thick] (-0.2,0) -- (8.4,0) node[right]{$t$};

  % --- gray CPs ---
  \foreach \x in {0.0, 0.5, 1.0, 1.5, 2.0, 2.5, 3.0, 3.5, 4.0,
                  4.5, 5.0, 5.5, 6.0, 6.5, 7.0, 7.5}
    \node[cp, fill=gray!30] at (\x, 0) {};

  % --- red (marginalized) CPs ---
  \foreach \x in {1.5, 2.0, 2.5}
    \node[cp, fill=red!55] at (\x, 0) {};

  % --- green (boundary) CPs ---
  \foreach \x in {3.0, 3.5, 4.0, 4.5}
    \node[cp, fill=green!55!black] at (\x, 0) {};

  % === ABOVE the timeline: window span brackets ===

  % Window w bracket: ticks at y=0.32, horizontal bar at y=0.55, label at y=0.75
  \draw[blue!75!black, thick]
    (1.5,0.32) -- (1.5,0.55) -- (4.5,0.55) -- (4.5,0.32);
  \node[blue!75!black] at (3.0, 0.80) {window $w$ \ (3\,s)};

  % Window w+1 bracket: ticks at y=0.95, bar at y=1.18, label at y=1.38
  \draw[blue!45, thick, dashed]
    (2.7,0.95) -- (2.7,1.18) -- (5.7,1.18) -- (5.7,0.95);
  \node[blue!45] at (4.2, 1.43) {window $w{+}1$};

  % === BELOW the timeline: zone labels ===
  % Red anchored east (text grows left); green anchored west (text grows right)
  % → labels always point away from the shared boundary at x=2.7

  \node[red!75!black, anchor=east] at (2.5, -0.65) {stride (0.3\,s)};
  \node[red!75!black, anchor=east, font=\footnotesize] at (2.5, -1.00) {marginalized};

  \node[green!55!black, anchor=west] at (2.9, -0.65) {boundary state};
  \node[green!55!black, anchor=west, font=\footnotesize] at (2.9, -1.00) {30-dim prior};

  % Prior propagation: boundary CPs feed the next window
  \draw[-{Latex[length=2mm]}, green!55!black, thick]
    (4.2, -1.30) to[out=0, in=180] (5.5, -1.30);
  \node[green!55!black, font=\footnotesize, anchor=west] at (5.55, -1.30)
    {prior $\to w{+}1$};

\end{tikzpicture}%
}% end scalebox
\caption{Sliding window: red CPs (stride zone) are marginalized
into a 30-dim Gaussian prior on the green boundary CPs.}
\label{fig:sw_diagram}
\end{figure}

The fixed-lag smoother advances a 3\,s window in 0.3\,s strides.
After each Ceres LM solve, the \emph{stride zone} CPs and orientation
knots are marginalized:
\begin{equation}
  S = H_{bb} - H_{ab}^{\top} H_{aa}^{-1} H_{ab},
  \label{eq:schur}
\end{equation}
where $a$ indexes stride-zone variables and $b$ indexes boundary
variables (last $N{-}1$ position CPs, last $N{-}1$ orientation
knots, and the 6-DoF bias; dimension 30).  $S$ is factored as
$S = LL^{\top}$ and stored as a \texttt{MargPriorFunctor} for the
next window.

\textbf{Re-centered prior.}
Before adding the prior to each window's problem, the linearization
point $\mathbf{x}_0$ is updated to the current warm-start.  This
makes the prior contribute only curvature (the Hessian shape), not a
gradient pull toward a stale historical estimate.  The residual is
$\mathbf{r} = L^{\top}(\mathbf{x} - \mathbf{x}_0)$ in local
coordinates.

\textbf{Consistent factor selection (Markov-blanket rule).}
Which residuals enter Eq.~\eqref{eq:schur} determines whether the
prior is a true marginal.  A na\"ive choice (all residuals touching
marginalized \emph{or} boundary variables) pulls in every \ac{imu} factor
of the window through the shared bias block ($\sim$25{,}000 residual
rows), \emph{conditions} on interior states (treating them as perfectly
known), and double-counts factors that are re-added in the next window.
The resulting prior is overconfident by orders of magnitude and must be
deflated by a hand-tuned scale of $10^{-7}$--$2{\times}10^{-4}$
depending on flight dynamics (\autoref{sec:prior_scale}).
The shared global bias is the trap: it is what tempts pulling in the
far-edge \ac{imu} factors.  We instead exploit the splines' local support
of $N$ knots (quintic position $N{=}6$, cubic orientation $N{=}4$).

\begin{proposition}[Exactly closed spline marginal]
\label{prop:marg}
Let $\mathcal{M}$ be the block of knots leaving the window and
$\mathcal{B}$ the $N{-}1$ retained knots immediately adjacent to it.
For an order-$N$ spline, every factor has local support over $N$
consecutive knots $[i,\,i{+}N{-}1]$.  If a factor touches $\mathcal{M}$,
i.e.\ $[i,\,i{+}N{-}1]\cap\mathcal{M}\neq\varnothing$, then
$[i,\,i{+}N{-}1]\subseteq\mathcal{M}\cup\mathcal{B}$.  Selecting exactly
the factors that touch $\mathcal{M}$, together with the previous prior,
therefore forms a Schur complement that conditions on no interior state
and shares no factor with the next window: the resulting prior is the
exact marginal.
\end{proposition}

Including \emph{only} residuals that touch a marginalized variable
(rather than those touching marginalized \emph{or} boundary variables)
thus yields a consistent marginal.  With this rule the prior is applied
\textbf{unscaled} (scale $=1$) across all flight regimes, and its
Mahalanobis residual at the solution drops from $10^{6}$--$10^{12}$
to $\mathcal{O}(1)$: the prior agrees with the incoming data.

\emph{Relation to \ac{fej}.}  This mechanism is distinct from, and
complementary to, first-estimates-Jacobian approaches to
marginalization consistency~\cite{dongsi2012consistency}: \ac{fej} fixes
the \emph{linearization points} of the prior to preserve the
observability nullspace, whereas our rule corrects the
\emph{factor set} entering the Schur complement so that the marginal
is exactly closed \emph{at a fixed linearization point}.  We do not
employ \ac{fej}; instead the prior is re-centered at each window's warm
start (curvature retained, gradient pull discarded).  We do not claim
this re-centering preserves consistency across relinearizations as a
theoretical guarantee; rather, in our experiments it left no measurable
inconsistency symptom; the \ac{nees} analysis of
\autoref{sec:nees} provides the corresponding empirical evidence.
Our missions are short (18--26\,s), which may be too brief to expose
drift a re-centered prior could accumulate over a long deployment;
bounding this over extended runs is left to future work.

\textbf{Live-edge warm start.}
Control points entering the window each stride would otherwise hold
stale dead-reckoned initialization values; any CP-level seam is
amplified by the min-snap term ($\propto \Delta/\Delta t_p^4$, initial
cost $\sim$$10^{13}$).  We rigidly align the entering segment (shift +
$\Delta R$) to the last data-supported CP of the previous solution,
reducing the initial cost by five orders of magnitude.

\textbf{Marginalization cost.}
The restricted Gauss--Newton Hessian for Eq.~\eqref{eq:schur} is
accumulated by evaluating the affected cost functions directly
(tangent-space Jacobians via the manifold plus-Jacobian, Triggs
correction for robust losses), avoiding a full solver re-linearization;
the marginalization step costs $\sim$10\,ms per window.
